Toen de jager introk
in een zelfgebouwd hok,

checkte hij geduldig
en vermaakte zorgvuldig

de nabije struiken.
Kon men onraad ruiken,

dan zat hij lang voor niets.
Hij had uitzicht op iets,

terwijl hij waterde
en zijn plas klaterde.

Deze ochtenduren
ging hij het begluren.

In zijn paradijsje
keek hij naar een meisje

dat haar vodden uittrok.
Zijn heimelijke hok

had ze niet opgemerkt.
“De camouflage werkt!”.

Ze was zo schaars gekleed
dat de jager niet meed

om de rest, ongezien,
van het meisje te zien.

Volgens wat in die tijd
gold als fatsoenlijkheid,

in de stedelijke kring,
waar de jager rondhing,

was ze immers reeds naakt,
alvorens kwijtgeraakt

haar geringe kleding
met schrale bedekking.

Ze droeg al nul dit wief
ofwel “Niks Um’t Lief”.

Dus werd Hume geboren,
opdat men kon horen

dat de zedelijkheid
geen objectiviteit,

maar emotioneel is.
De jager had het mis.

Zijn normatief oordeel
was slechts contextueel.

Deze roodharige
net meerderjarige

was schijnbaar schaamteloos.
Er was naar haar iets loos.

Schoon ze haar kledij stak
aan een gebroken tak

van een rijzige den,
voelde ze zich niet zen.
Nog bij de dennenboom
toonde ze enige schroom

om zich te bewegen.
Enigszins verlegen

verliet het wicht de den.
Ze bewoog naar een ven.

Ondanks de felle kou
was het klaar wat ze wou.

Het oeverveen deinde
bij elk voeteneinde

dat het zachtjes roerde
alsof het vervoerde

een verheven wezen,
die ook werd geprezen

door de mist die optrok
om voor het jagershok

haar schoonheid te klaren
en te openbaren.

Bij elke stap nader
bleek ze een aanrader.

Deze parel, zo puur
gevormd door de natuur,

was om naar te staren.
De kleur van haar haren
wist te benaderen
die van herfstbladeren.

Haar huid had veel sproetjes.
Haar handjes en voetjes

zwommen als heikneuter,
spartelend als kleuter,

zodat de roodharige
haar dertienjarige

ogende bips uitstak
boven het oppervlak

van het vennewater.
Het wicht draaide later

haar billen vliegensvlug
op haar ijskoude rug,

zodat zo haar voorkant
uit het water beland.

Het meest haar twee kuppels,
die als halve druppels

met ieder een mesa,
zoals uit Nevada,

fel rood gekleurd en stijf,
en ver van het voorlijf,

als kaarsrechte staken
de moerasmist raakten.

Zonder verrekijker
glunderde de kijker

toen ze van hem afzwom
en haar marmotje glom.

Uit zijn linnen knapzak
haalde hij een flink pak

papieren tevoorschijn.
Het openbaar bloot-zijn

en de gouden ketting,
die los om haar hals hing

met daaraan twee kralen
met de initialen

van Emmy Nativus
naast het kruis van Christus,

maakte het herkennen
moeilijk te ontkennen.

Ze voldeed aan de claims.
Ze was beslist inheems.

Al was ze een Rotschopf,
met het haar in een Zopf,

toch galmde haar klankmat
schel het Niederdeutsch Platt,

toen ze op speels gemak
tegen een eendje sprak.

Hij kon dit gouden kalf,
al zag hij het maar half,

volledig aanbidden.
“Focus op het midden!”

De weg tussen genot
en nageleefd verbod

bracht hem vast verlichting.
Luister als zingeving.

De driehoekige mat,
die op haar schaamstreek zat,

werd gevormd door haartjes
in krullende staartjes

met een fel oranje tint.
Ze zwiepten door de wind

en werden uitgestrekt.
Ook haar kruin werd bedekt

door rode haarlokken,
die de aandacht trokken

en hingen langs haar wang
rond een armlengte lang.

Haar gloed was innemend
en adembenemend.

Nimmer zag hij een wicht
met zo’n lieflijk gezicht.

Toen een ijzig zuchtje
in een dwarrelvluchtje

over haar nippeltjes,
omringd door stippeltjes,

haar strelend kietelde,
als vrouw betitelde,

dwaalden haar gedachten
naar vergeten smachten.

Naar dezelfde prikkel,
die een stoere bikkel

liefdevol gaf aan haar.
Al strelend met een aar.

De lange stengel gras
maakte een zachte kras

op Emmy’s tepelhof.
Waar de kolf haar huid trof.

In de kilte waden
gevolgd door windbaden

bracht haar intens genot,
maar verduisterde God.

Welke hersenkogels
schoten de roofvogels,

die sereen cirkelden
over verzinselen?

“Kafka-personage”
was nog geen passage

in haar liefdesverhaal.
Zou ze na dit etmaal

wel zijn veroordeeld,
zonder te zijn gedeeld

in wie de stille kracht
was achter de jacht?
